\documentclass{article}
\usepackage[style=apa]{biblatex}
\usepackage{graphicx}
\usepackage{subcaption}
\usepackage{float}
\usepackage{amsmath}
\addbibresource{references.bib}

\title{CS6480 Final Report}
\author{Talmage Gaisford}
\date{April 2026}

\begin{document}

\maketitle

\section{Goals}

Pellet grills present a difficult control problem due to delayed 
response times, variable combustion, and disturbances such as lid openings.
This project evaluates whether a neural network–based controller can
improve temperature regulation compared to a classical PID controller.

A physics-based simulation of a pellet grill was developed to test both
approaches. A Soft Actor-Critic (SAC) agent was trained to control fan
speed and pellet feed, and its performance was compared to a tuned PID
controller across several cooking scenarios using offset from the target
temperature as the primary metric.


\section{Background}

Pellet grills, such as the Traeger Ironwood 650, regulate temperature by
controlling airflow and fuel delivery, typically using PID-based control.
Although PID controllers are simple, effective, and well tested, they can struggle with 
disturbances present in the system and can cause temperature oscillation.

Disturbances like lid openings cause rapid heat loss, and PID
controllers often react too slowly or overshoot during recovery. Manufacturers are opposed
to adding additional sensors for the sake of keeping points of failure down; improving control
performance through software alone is desirable.

Reinforcement learning provides an alternative approach by learning control
policies directly from system behavior. This project applies a Soft
Actor-Critic (SAC) agent to determine whether a learned controller can
better handle these challenges than a traditional PID controller.


\section{Methods}

\subsection{Simulation Model}

A physics-based pellet grill simulator was developed to evaluate two
temperature control strategies.

Fan speed (0--1) controls airflow into the firepot and follows a
first-order response with a 5-second time constant. The auger is binary
(on/off) and delivers fuel with a 3-second response delay.

Combustion rate is limited by the minimum of airflow and fuel supply,
reflecting the requirement for both inputs. Heat output is proportional
to combustion, with a maximum of 45 heat units.

Heat loss follows Newton's law of cooling (coefficient 0.1 per
\textdegree F above ambient). Lid openings increase heat loss by $4\times$,
and fan speed adds a small convective loss term (0.01).

Grill temperature evolves from net heat input divided by a thermal mass
of 50~\textdegree F{\textperiodcentered}s per unit. The simulator runs at
1-second timesteps, with varying ambient conditions and lid events.

\subsection{PID Controller}

A PID controller with anti-windup was implemented and tuned via parameter
sweep on a 60-minute, 225\textdegree F reference cook. Final gains were
$k_p = 0.015$, $k_i = 0.002$, and $k_d = 0.2$.

The controller outputs a single signal in $[0,1]$ that directly sets fan
speed. The auger switches on when this signal exceeds 0.5, coupling airflow
and fuel control.

\subsection{Neural Network Controller (SAC)}

A Soft Actor-Critic (SAC) agent (using \texttt{stable\_baselines3}) was
trained to control the system. The agent observes:

\begin{equation*}
    \bigl[T_{\text{grill}},\;\dot{T},\;\ddot{T},\;
    v_{\text{fan}},\;u_{\text{auger}},\;
    T_{\text{target}},\;e,\;\tau\bigr]
\end{equation*}

where $e = T_{\text{target}} - T_{\text{grill}}$ and $\tau$ is normalized
episode time. Ambient temperature and lid state are not observed.

The action space is continuous: $[a_{\text{fan}}, a_{\text{auger}}] \in [0,1]^2$.
Fan maps directly to speed, while auger output is thresholded at 0.5 to
produce binary behavior.

The reward is:
\begin{equation*}
    r_t = -\left| T_{\text{grill}} - T_{\text{target}} \right|
\end{equation*}

Training used 1,000,000 timesteps, learning rate $10^{-4}$, 3,600-step
episodes, and 4 parallel environments. The best model was selected based
on evaluation performance every 10,000 steps.

\subsection{Evaluation Scenarios}

Four scenarios were evaluated:

\begin{enumerate}
    \item \textbf{Steaks} --- 30~min, 450\textdegree F, 65\textdegree F ambient, one lid event
    \item \textbf{Chicken} --- 90~min, 375\textdegree F, 75\textdegree F ambient, lid every 30~min
    \item \textbf{Ribs} --- 3~hours, 275\textdegree F, 40\textdegree F ambient, spritz at 1h and 2h
    \item \textbf{Brisket} --- 6~hours, 225\textdegree F, 25\textdegree F ambient, lid every 30~min
\end{enumerate}

Each run includes a preheat phase (excluded from evaluation) followed by
the cook phase. Performance is measured using mean absolute error (MAE).

\section{Results}

Representative temperature traces for each scenario are shown in
Figure~\ref{fig:all_scenarios}.

\begin{figure}[H]
    \centering

    \begin{subfigure}{0.48\textwidth}
        \centering
        \includegraphics[width=\linewidth]{brisket.png}
        \caption{Brisket (6h, 225\textdegree F)}
    \end{subfigure}
    \hfill
    \begin{subfigure}{0.48\textwidth}
        \centering
        \includegraphics[width=\linewidth]{chicken.png}
        \caption{Chicken (90 min, 375\textdegree F)}
    \end{subfigure}

    \vspace{0.5em}

    \begin{subfigure}{0.48\textwidth}
        \centering
        \includegraphics[width=\linewidth]{ribs.png}
        \caption{Ribs (3h, 275\textdegree F)}
    \end{subfigure}
    \hfill
    \begin{subfigure}{0.48\textwidth}
        \centering
        \includegraphics[width=\linewidth]{steaks.png}
        \caption{Steaks (30 min, 450\textdegree F)}
    \end{subfigure}

    \caption{Temperature control comparison between PID and neural network controllers across all scenarios.}
    \label{fig:all_scenarios}
\end{figure}



Table~\ref{tab:mae} reports the mean absolute temperature error for each
controller across all four scenarios.

\begin{table}[h]
    \centering
    \begin{tabular}{lllrrl}
        \hline
        \textbf{Scenario} & \textbf{Duration} & \textbf{Target} &
        \textbf{PID MAE (\textdegree F)} & \textbf{NN MAE (\textdegree F)} &
        \textbf{Outcome} \\
        \hline
        Steaks  & 30~min   & 450\textdegree F & 8.50 & 10.10 & PID wins \\
        Chicken & 90~min   & 375\textdegree F & 2.10 &  1.90 & NN wins  \\
        Ribs    & 3~hours  & 275\textdegree F & 1.09 &  1.13 & $\approx$ Tie \\
        Brisket & 6~hours  & 225\textdegree F & 0.50 &  0.30 & NN wins  \\
        \hline
    \end{tabular}
    \caption{Mean absolute temperature error (MAE) by scenario. Lower is better.}
    \label{tab:mae}
\end{table}

The NN controller has the advantage on longer cooks. The
performance gap widens with cook duration: the NN edges out the PID on
chicken, matches it on ribs, and nearly halves PID error on the 6-hour
brisket cook.

Figure~\ref{fig:all_scenarios} shows temperature performance, but it does
not reveal how each controller achieves that behavior. To better understand
how control is achieved, Figure~\ref{fig:ribs_controls} shows the internal
control signals for the ribs scenario.

The PID controller produces a single coupled control signal, resulting in
synchronized adjustments to fan speed and auger activity. In contrast, the
neural network independently modulates airflow and fuel delivery, allowing
it to apply short bursts of fuel while maintaining smoother fan control.
This separation enables more precise handling of disturbances such as lid
openings.

\begin{figure}[H]
    \centering
    \includegraphics[width=\linewidth]{ribs_controls.png}
    \caption{Controller behavior during the ribs scenario (3~hours, 275\textdegree F).
    Top: grill temperature response. Middle: fan speed. Bottom: auger duty cycle.
    The neural network (SAC) independently adjusts fan and auger signals,
    while the PID controller produces coupled responses.}
    \label{fig:ribs_controls}
\end{figure}

\section{Conclusions}

The neural network controller outperforms the PID on longcooks. 
Its advantage comes partially from its ability to independently control fan
speed and auger feed, something the coupled PID cannot do. The NN
learned to manage the lag between airflow and fuel delivery and to anticipate
disturbances (ambient temperature, lid events) from temperature derivatives
alone.

The PID remains competitive for medium-length cooks and is simpler to
implement. With properly tuned parameters it achieved good
accuracy on the ribs and chicken scenarios.

On short, high-temperature cooks (steaks, 30~minutes, 450\textdegree F),
the PID's reactive control beats out the NN. The NN was designed and trained to manage sustained,
steady low-and-slow cooks where anticipating and smoothing slow thermal drift
matters. At high temperatures near the grill's maximum, the
control problem reduces to a simpler task that a well-tuned reactive
controller handles efficiently.

\end{document}
